% La metrique : produit scalaire sur l'espace tangent, donc norme et angle,
% et par integration longueur d'une courbe.
% Utilise dans la \cref{sec:context-riemann-notion}.
% Styles et briques communs aux figures de geometrie riemannienne du
% \cref{ch:context-riemann}. Ce fichier peut etre \input plusieurs fois.
\definecolor{rmanifold}{RGB}{90,90,90}
\definecolor{rtangent}{RGB}{50,110,120}
\definecolor{rgeodesic}{RGB}{40,70,140}
\definecolor{rvector}{RGB}{175,95,45}
\tikzset{
  riemann/.style={
    >={Stealth[length=4.5pt,width=3.5pt]},
    line join=round,
    font=\footnotesize,
  },
  % La nappe qui represente la variete, centree sur l'origine.
  rsurface/.pic={
    \fill[rmanifold!7]
      (-3.2,0.15) .. controls (-2.95,1.75) and (-1.5,2.35) .. (0,2.35)
      .. controls (1.5,2.35) and (2.95,1.7) .. (3.2,0.15)
      .. controls (1.9,-0.7) and (-1.9,-0.7) .. (-3.2,0.15) -- cycle;
    \draw[rmanifold!70, thick]
      (-3.2,0.15) .. controls (-2.95,1.75) and (-1.5,2.35) .. (0,2.35)
      .. controls (1.5,2.35) and (2.95,1.7) .. (3.2,0.15);
    \draw[rmanifold!70, thick]
      (-3.2,0.15) .. controls (-1.9,-0.7) and (1.9,-0.7) .. (3.2,0.15);
  },
  % Le plan tangent, a placer en un point donne via un shift.
  rplane/.pic={
    \fill[rtangent!10, draw=rtangent!65]
      (-1.35,-0.42) -- (1.55,-0.88) -- (2.05,0.22) -- (-0.85,0.68) -- cycle;
  },
}

\begin{tikzpicture}[riemann]
  \coordinate (P) at (-1.70,1.15);

  \path (0,0) pic {rsurface};
  \node[rmanifold!85] at (2.85,1.15) {$\Manifold$};

  % plan tangent et deux vecteurs
  \path (P) pic[scale=0.75] {rplane};
  \draw[rvector, very thick, ->] (P) -- ($(P)+(1.10,-0.16)$);
  \node[rvector, anchor=north west, inner sep=1pt] at ($(P)+(0.85,-0.14)$) {$\tvec$};
  \draw[rvector, very thick, ->] (P) -- ($(P)+(0.40,0.46)$);
  \node[rvector, anchor=south east, inner sep=1pt] at ($(P)+(0.42,0.44)$) {$\tvecbis$};
  \draw[rtangent!80] ($(P)+(0.46,-0.07)$) arc[start angle=-8, end angle=48, radius=0.47];
  \node[rtangent!90, font=\scriptsize, anchor=west] at ($(P)+(0.80,0.42)$)
    {$\metric{\tvec}{\tvecbis}{\param}$};
  \node[rtangent!90, anchor=north, inner sep=2pt] at ($(P)+(0.10,-0.50)$)
    {$\Tangent{\param}{\Manifold}$};

  % longueur d'une courbe : vitesses echantillonnees le long du chemin
  \draw[rgeodesic, very thick]
    (0.55,1.85) .. controls (1.45,1.95) and (2.05,1.55) .. (2.45,0.95);
  \node[rgeodesic, anchor=south east, inner sep=1pt] at (0.60,1.88) {$\geod$};
  \foreach \x/\y/\dx/\dy in {0.85/1.91/0.30/0.05, 1.60/1.89/0.28/-0.13, 2.22/1.42/0.18/-0.28}
    \draw[rvector!85, thick, ->] (\x,\y) -- ++(\dx,\dy);
  \node[rgeodesic!85, anchor=north, font=\scriptsize] at (1.35,1.32)
    {$L(\geod)=\int\normriem{\dot{\geod}}{\geod}$};

  \fill (P) circle (1.7pt);
  \node[anchor=south east, inner sep=2pt] at (P) {$\param$};
\end{tikzpicture}
